% ===================================================================
%  اللوحة نمرة ١١ — المدينة ومبانيها. --- الحريقة.
%
%  Traduction, faite en 2026, en arabe egyptien ecrit, sur l'ido de
%  texto/io. Regles de la colonne : voir 10-lawha-01.tex. Deux
%  parties, deux numerotations : les cles portent c1 et c2.
%
%  LA VILLE EST L'ENDROIT OU L'EGYPTE A SES PROPRES ADMINISTRATIONS,
%  ET ELLES TOMBENT UNE A UNE SUR CELLES DU LIVRET DE 1926 :
%
%    مديرية (15)     la prefecture — et le prefet (91) est المدير
%    مأمور (94)      le commissaire de police
%    وكيل النيابة (95) le magistrat qui instruit
%    بوليس (64)      les agents        MSA : شرطة
%    قشلاق (16)      la caserne        — du turc kışla
%    بوسطجي (44)     le facteur        — le suffixe -جي, encore
%    بوستة (43)      la poste
%    دار المعلّمين (41) l'ecole normale
%    بورصة (14)      la Bourse
%
%  Ce n'est pas une correspondance approchee. La مديرية egyptienne
%  EST la circonscription administrative d'un chef-lieu de
%  departement, et le مدير qui la dirige est nomme comme le prefet
%  l'etait ; le مأمور est le chef de la police d'un district ; وكيل
%  النيابة est celui qui ouvre l'enquete. La hierarchie du dernier
%  bloc se lit donc sans une glose.
%
%  ET LA VILLE ORDINAIRE A SES MOTS :
%
%    ميدان (96)   la place
%    حارة (11)    la ruelle
%    قهوة (25)    le cafe — en Egypte, le lieu s'appelle du nom de
%                   la boisson
%    تياترو (73)  le theatre    de l'italien teatro
%    حريقة (74)   l'incendie
%    مطافي (86)   les pompiers
%    كشك (27, 55) le kiosque
%    عفش (87)     les meubles
%    كمنجة (81)   le violon
%
%  C'EST LE PENDANT DU TABLEAU 11 CANTONAIS, ou 差人, 警司 et 郵差
%  faisaient le meme travail. Deux villes, deux administrations, et
%  chaque fois la meme demonstration : une langue qui gouverne a des
%  noms pour ses fonctions.
%
%  LE FAC-SIMILE DONNE DEUX FOIS LE RENVOI (81) au bloc c2-06-1 — aux
%  musiciens, puis au violon. La traduction le rend deux fois : elle
%  vise les memes objets dans le meme ordre.
% ===================================================================

%%K t11-apar-1 apar
\VUsaut{2.0mm}
\VUcentre{13.2pt}{90}{السلسلة التالتة}
\VUsaut{3.0mm}
\VUfilet{16mm}
\VUsaut{5.6mm}
\VUtitre{15.0pt}{60}{اللوحة نمرة ١١}
\VUsaut{1.4mm}
\VUpk{11.4pt}{40}{المدينة ومبانيها. --- الحريقة.}
\VUsaut{2.2mm}

%%K t11-tit-1 sub
\VUpk{10.2pt}{40}{الضاحية.}
\VUsaut{3.0mm}

%%K t11-c1-01-1 p
1. --- أنا ساكن في مدينة كبيرة على بعد مية كيلومتر من المحيط، ومع كده هي
\VUgras{ميناء} \textsuperscript{(1)} من الدرجة الأولى. والنهر اللي على حدّها عرضه
أربعميت متر وبيعمل مرسى حلو أوي.

%%K t11-c1-02-1 p
2. --- والجزء كله من المدينة اللي على \VUgras{الأرصفة} \textsuperscript{(2)} شكله كله
بحري وصناعي، وبيعمل \VUgras{ضاحية} \textsuperscript{(3)} ما بيسكنهاش غير البحّارة
والعمال. ودول بيشتغلوا في \VUgras{المصانع} \textsuperscript{(4)} وفي \VUgras{مصنع
الغاز} \textsuperscript{(5)} اللي \VUgras{مدخنته} \textsuperscript{(6)} العالية
و\VUgras{خزان الغاز} بتوعه بيتشافوا من بعيد أوي.

%%K t11-c1-03-1 p
3. --- وفي العصور الوسطى، كانت المدينة متحصّنة، ولسه بيتلاقى كام أثر من
أسوارها، وبالذات \VUgras{باب} \textsuperscript{(8)} \VUgras{القلعة} \textsuperscript{(7)}
القديمة اللي على جنبيه \VUgras{برجين} \textsuperscript{(9)} بيستعملوهم دلوقت
\VUgras{سجن}.

%%K t11-c1-04-1 p
4. --- وفي \VUgras{الحيّ} ده كله، اللي شكله قديم أوي، مافيش غير مبنى واحد جديد
نسبيًا: وهو \VUgras{الجمرك} \textsuperscript{(10)}. وهو بين الرصيف و\VUgras{الحارة}
\textsuperscript{(11)} اللي بتودّي لـ\VUgras{دار البلدية} \textsuperscript{(12)} القديمة.

%%K t11-c1-04-1 p suite
والمبنى ده سهل تعرفه من \VUgras{برج الساعة} \textsuperscript{(13)} العملاق اللي
بيشرف على المدينة القديمة.

%%K t11-c1-05-1 p
5. --- وعلى الناحية التانية من الشارع، \VUgras{واقف} مبنى مبني بنفس طراز الجمرك:
وده \VUgras{البورصة} \textsuperscript{(14)}. وواحدة من واجهاته بتطلّ على ميدان صغير
فيه \VUgras{المديرية} \textsuperscript{(15)}. والحقيقة إن مدينتنا، من غير ما تبقى
عاصمة كبيرة، هي مع كده عاصمة مديرية مهمة.

%%K t11-tit-2 sub
\VUsaut{5.0mm}
\VUpk{10.2pt}{40}{أعلى جزء من المدينة.}
\VUsaut{3.4mm}

%%K t11-c1-06-1 p
6. --- والحامية، وعددها مش قليل، ساكنة في \VUgras{قشلاقات} \textsuperscript{(16)}
واسعة وصحّية أوي؛ وهي بتمتدّ لحدّ سور \VUgras{الجنينة العامة} \textsuperscript{(17)}.
و\VUgras{الشارع الكبير} \textsuperscript{(19)} اللي بيودّي للجنينة دي بيعدّي ورا
\VUgras{صندوق التوفير} \textsuperscript{(18)} وبينتهي عند ميدان \VUgras{الكاتدرائية}
\textsuperscript{(20)}.

%%K t11-c1-07-1 p
7. --- والمباني دي كلها بتمتدّ زي المدرّج على تلّ صغير، اللي فيه كمان
\VUgras{مدرستنا الثانوية} \textsuperscript{(21)} الواسعة.

%%K t11-c1-07-2 p
ولو نزلت من طريق مايل شوية، بيوصل للشارع الكبير، توصل لـ\VUgras{المحكمة}
\textsuperscript{(22)} اللي قدامها بتمتدّ \VUgras{جنينة عامة} \textsuperscript{(26)}
مريحة. و\VUgras{الجنينة} دي مش كبيرة أوي، بس بتعمل في وسط المدينة واحة خضرا
وبرد. علشان كده أهل المدينة بيترددوا عليها أوي، وبيحبّوا ييجوا يقعدوا تحت خيمة
\VUgras{القهوة} \textsuperscript{(25)}، علشان يسمعوا عازفين العسكرية اللي بيعزفوا يوم
الخميس ويوم الحد في \VUgras{الكشك} \textsuperscript{(27)} اللي بين النجيلة والأحواض.

%%K t11-c1-08-1 p
8. --- وعلى واحدة من النجيلات دي واقف \VUgras{تمثال} \textsuperscript{(28)} برونز،
اتعمل لذكرى واحد من أهل بلدنا، اللي عمل خدمات كبيرة لبلده اللي اتولد فيها.

%%K t11-tit-3 sub
\VUsaut{4.0mm}
\VUpk{10.2pt}{40}{المواصلات.}
\VUsaut{3.4mm}

%%K t11-c1-09-1 p
9. --- ولحدّ دلوقت مدينتنا لسه مش متنوّرة بالكهربا، عندها بس \VUgras{لمبات الغاز}
\textsuperscript{(29)}. إنما \VUgras{الجرايد المحلّية} بتعمل حملة علشان طريقة
الإنارة دي تتحسّن، \VUgras{زي} ما اتحسّنت \VUgras{طريقة جرّ الترام}.
و\VUgras{العربيات} القديمة، اللي كانت \VUgras{بتجرّها} \VUgras{الخيل}، اتغيّرت
عمليًا بـ\VUgras{الترام} \textsuperscript{(31)} اللي بينقل الركّاب بسرعة من طرف
\VUgras{المدينة للطرف التاني}، بسعر \VUgras{رخيص} أوي.

%%K t11-c1-10-1 p
10. --- وجنب دار القضا في \VUgras{البنك} \textsuperscript{(32)} اللي بيخصموا فيه
الأوراق التجارية. و\VUgras{محصّلين البنك} \textsuperscript{(33)} بيروحوا يحصّلوا
الديون في البيوت.

%%K t11-tit-4 sub
\VUsaut{4.0mm}
\VUpk{10.2pt}{40}{الفن والتعليم.}
\VUsaut{3.4mm}

%%K t11-c1-11-1 p
11. --- والمبنى الحلو اللي واقف قريّب هو \VUgras{المتحف}. وهو فيه في نفس الوقت
\VUgras{مكتبة} \textsuperscript{(34)} المدينة، و\VUgras{مجموعات التاريخ الطبيعي}
\textsuperscript{(35)} الحلوة (متحف الطبيعة)، و\VUgras{متحف صور} \textsuperscript{(36)}
لطيف اللي هو \VUgras{معرض} دايم رائع.

%%K t11-c1-11-2 p
وبيفتح للجمهور مرتين في الأسبوع، إنما الأجانب يقدروا يزوروه أي يوم مقابل بقشيش
لـ\VUgras{الغفير} \textsuperscript{(37)}. و\VUgras{الرسّامين} \textsuperscript{(38)}
بييجوا هناك يشتغلوا. وبتشوفهم \VUgras{قدام} حمّالاتهم، و\VUgras{البالتّة} في
إيدهم، بينقلوا اللوحات المشهورة.

%%K t11-c1-12-1 p
12. --- وعلى المرتفع، ورا المتحف، بتبان \VUgras{قبّة} \textsuperscript{(40)}
\VUgras{المستشفى} \textsuperscript{(39)} ومباني \VUgras{دار المعلّمين}
\textsuperscript{(41)} اللي اتدرّب فيها مدرّسين مدارسنا البلدية. وفي ميدان التياترو
في \VUgras{بوستة} \textsuperscript{(43)} متركّبة في \VUgras{بيت خاص}
\textsuperscript{(42)}.

%%K t11-tit-5 sub
\VUsaut{5.0mm}
\VUpk{11.4pt}{40}{الحريقة.}
\VUsaut{3.0mm}

%%K t11-tit-6 sub
\VUpk{10.2pt}{40}{«حريقة!»}
\VUsaut{3.4mm}

%%K t11-c2-01-1 p
1. --- خبر مرعب اتنشر في المدينة: حريقة حصلت حالًا في التياترو! وفي اللحظة اللي
وصلت فيها \VUgras{الميدان} \textsuperscript{(96)}، الزحمة كانت اتجمّعت خلاص على
\VUgras{الرصيف} \textsuperscript{(46)} وعلى \VUgras{الشارع} \textsuperscript{(47)}.
و\VUgras{البوسطجية} \textsuperscript{(44)} و\VUgras{سعاة الخطابات} \textsuperscript{(45)}
خرجوا من البوستة. و\VUgras{سوّاق الترام} \textsuperscript{(48)} اضطرّ يوقّف عربيته
علشان يعدّي \VUgras{عربية الإسعاف} \textsuperscript{(50)} البلدية، اللي جاية
بـ\VUgras{جرّاح} \textsuperscript{(52)} و\VUgras{ممرّض} \textsuperscript{(51)} علشان
يعالجوا \VUgras{مصاب} \textsuperscript{(53)} اتنقل على \VUgras{نقّالة}
\textsuperscript{(54)} جنب \VUgras{كشك الجرايد} \textsuperscript{(55)}.

%%K t11-c2-02-1 p
2. --- وسريّة مشاة كانت راجعة من التمرين وقفت علشان تأمّن النظام مع
\VUgras{البوليس الخيّال} \textsuperscript{(61)} البلدي. و\VUgras{الخيّالة} دول، اللي
خيلهم بتقف على رجلين، أحسن من \VUgras{العساكر} \textsuperscript{(57)} أو
\VUgras{الجندرمة} \textsuperscript{(63)} في إنهم يرجّعوا \VUgras{الفضوليين}
\textsuperscript{(56)} لورا.
وكمان الجندرمة مشغولين أوي. وواحد فيهم بيوضّح لـ\VUgras{البوليس}
\textsuperscript{(64)} فين لازم ينقلوا \VUgras{مصاب} \textsuperscript{(65)} تاني، واحد
من رجّالة المطافي الشجعان وقع من على سلّم.

%%K t11-c2-03-1 p
3. --- و\VUgras{الضابط} \textsuperscript{(66)} بيحدّد المكان اللي لازم تتحطّ فيه
\VUgras{طرمبة البخار} \textsuperscript{(67)} الجاية. ولحدّ ما الماكينة القوية دي
توصل، خلّى الرجّالة يركّبوا \VUgras{طرمبة إيد} \textsuperscript{(68)}،
\VUgras{خرطومها} \textsuperscript{(69)} الطويل بيرمي المية زي السيل على النار.

%%K t11-c2-03-1 p suite
وستّة من المطافي بيشغّلوها، وزميلهم اللي ماسك \VUgras{فوهة الخرطوم}
\textsuperscript{(70)} بيوجّه \VUgras{عمود المية} \textsuperscript{(71)} لوسط الحريقة.

%%K t11-c2-04-1 p
4. --- وعلى الناحية التانية من التياترو، سندوا \VUgras{السلّم} \textsuperscript{(72)}
الكبير. وده بيخلّي المطافي يقدروا يقرّبوا من ألسنة النار اللي بتطلع وسط دوّامات
\VUgras{دخان} \textsuperscript{(75)} أسود.

%%K t11-tit-7 sub
\VUsaut{5.0mm}
\VUpk{10.2pt}{40}{أعمال بطولة.}
\VUsaut{3.4mm}

%%K t11-c2-05-1 p
5. --- \VUgras{الحريقة} \textsuperscript{(74)} ابتدت في صالة انتظار
\VUgras{التياترو} \textsuperscript{(73)}، وهما بيبروفوا \VUgras{الأوبرا} اللي
\VUgras{أول عرض} ليها كان هيبقى بكرة. ولحسن الحظ ما كانش في غير كام
\VUgras{متفرّج} \textsuperscript{(76)} في الصالة. والمساكين دول لجأوا لـ\VUgras{البلكونة}
\textsuperscript{(77)}، اللي جه \VUgras{رجّالة المطافي} \textsuperscript{(86)} الشجعان
ينجّدوهم فيها بالسلّم والحبال.

%%K t11-c2-06-1 p
6. --- وتحت \VUgras{الرواق} \textsuperscript{(78)}، اللي \VUgras{الإعلان}
\textsuperscript{(79)} بيقول فيه عن العرض، الواحد بيشوف \VUgras{عازفين}
\textsuperscript{(81)} \VUgras{الأوركسترا} \textsuperscript{(80)} بيهربوا وهما شايلين
آلاتهم: \VUgras{كمنجة} \textsuperscript{(81)}، و\VUgras{فلوت} \textsuperscript{(82)}،
و\VUgras{تشيلو} \textsuperscript{(83)}، و\VUgras{ترمبون} \textsuperscript{(84)}،
و\VUgras{طبلة} \textsuperscript{(85)}، وهكذا. والناس بتحاول تنقذ جزء من
\VUgras{العفش} \textsuperscript{(87)}.

%%K t11-c2-07-1 p
7. --- و\VUgras{الممثّلين} \textsuperscript{(89)} و\VUgras{الممثّلات}
\textsuperscript{(88)}، و\VUgras{المغنّيين} و\VUgras{المغنّيات}، اضطرّوا يهربوا
بـ\VUgras{كستيوم} \textsuperscript{(90)} المسرح. والتينور بيشرح لـ\VUgras{صحفي}
\textsuperscript{(92)} إزاي حصل ده. و\VUgras{المراسل} جري من أول لحظة مع أصحاب
السلطة: \VUgras{المدير} \textsuperscript{(91)}، و\VUgras{رئيس البلدية}
\textsuperscript{(93)}، و\VUgras{مأمور البوليس} \textsuperscript{(94)}، و\VUgras{وكيل
النيابة} \textsuperscript{(95)}، اللي هيحقّقوا علشان يحدّدوا المسؤوليات.

\VUsaut{6.0mm}
\VUfilet{22mm}
